\chapter{Tecnologías usadas}
\label{ch:tech}

Para implementar este tipo de sistemas requerimos de tecnologías que nos permitan cumplir con todos los requisitos especificados y el diseño mostrado en el capítulo anterior. Para cada componente veremos: una breve explicación de la necesidad de dicho componente, listado de las distintas alternativas ideales que existen para dicho componente, una tabla comparativa para distintos aspectos técnicos y la tecnología seleccionada con su correspondiente justificación. 

\section{Servidor back-end}

Para el servidor central que gestiona la lógica de negocio y expone la API REST, se requiere un sistema fuertemente tipado, escalable y con un ecosistema maduro de dependencias.

\subsection*{Evaluación de alternativas}

\subsubsection*{Spring Boot (Kotlin)}
Framework empresarial maduro que destaca por su robusto sistema de inyección de dependencias y módulos integrados (Spring Security, Data JPA). El uso del lenguaje Kotlin aporta un tipado estático moderno, inmutabilidad y seguridad contra referencias nulas (\emph{Null-safety}) de forma nativa.

\subsubsection*{Node.js (NestJS)}
Entorno de ejecución asíncrono basado en TypeScript. Ofrece tiempos de arranque muy rápidos y un ecosistema masivo de paquetes. Su arquitectura basada en eventos es ideal para alta concurrencia de entrada/salida, aunque delega gran parte de la estructura empresarial a librerías de terceros.

\subsubsection*{Ktor}
Framework asíncrono ligero creado por JetBrains para Kotlin. Ofrece un alto rendimiento mediante el uso de corrutinas y un diseño altamente modular, pero carece de las soluciones empresariales completas que proporciona el ecosistema Spring.

\subsection*{Comparativa técnica}

En la Tabla \ref{tab:comparativa-backend} se resumen las características fundamentales de cada entorno.

\begin{table}[H]
    \hspace*{-0.25cm}
    \begin{tabular}{|p{0.22\textwidth}|p{0.22\textwidth}|p{0.22\textwidth}|p{0.22\textwidth}|}
        \hline
        \diagbox[width=\dimexpr 0.22\textwidth+2\tabcolsep\relax]{\textbf{Criterio}}{\textbf{Tecnología}} & \textbf{Spring Boot (Kotlin)} & \textbf{NestJS (TypeScript)} & \textbf{Ktor (Kotlin)} \\ \hline
        \textbf{Paradigma} & Multihilo / Síncrono & Basado en eventos & Asíncrono (Corrutinas) \\ \hline
        \textbf{Seguridad de tipos} & Estricta (Nativa) & Alta (Transpilado) & Estricta (Nativa) \\ \hline
        \textbf{Ecosistema} & Muy maduro & Masivo (NPM) & En crecimiento \\ \hline
        \textbf{Arquitectura} & Altamente estructurada & Modular & Minimalista \\ \hline
    \end{tabular}
    \caption{Comparativa de frameworks para el desarrollo del back-end.}
    \label{tab:comparativa-backend}
\end{table}

\subsection*{Justificación de la elección}

Se ha seleccionado \textbf{Spring Boot con Kotlin} por los siguientes motivos técnicos:
\begin{itemize}
    \item \textbf{Diseño arquitectónico:} Su robusta inyección de dependencias facilita la implementación estricta de la Arquitectura Hexagonal, permitiendo un desacoplamiento real entre los puertos y sus adaptadores.
    \item \textbf{Mitigación de riesgos:} Contar con experiencia previa en el ecosistema Spring reduce los tiempos de configuración y asegura una implementación estable del sistema de seguridad e interacciones con la base de datos.
    \item \textbf{Modernización técnica:} La adopción de Kotlin representa un reto académico de alto valor, permitiendo escribir un código de dominio más conciso y expresivo frente a lenguajes tradicionales como Java.
\end{itemize}

\section{Base de datos}

Para dar soporte a la persistencia del sistema se requiere un Sistema de Gestión de Bases de Datos Relacionales (SGBDR) que garantice el cumplimiento de las propiedades \ac{ACID}.

\subsection*{Evaluación de alternativas}

\subsubsection*{PostgreSQL}
Sistema objeto-relacional de código abierto con adherencia estricta a los estándares SQL. Destaca por su alta extensibilidad, soportando tipos de datos avanzados y manipulación nativa de estructuras semiestructuradas mediante \texttt{JSONB}. Su modelo de control de concurrencia multiversión (MVCC) optimiza el rendimiento en entornos de alta carga transaccional.

\subsubsection*{MySQL}
Uno de los motores relacionales más extendidos, basado en el motor de almacenamiento \emph{InnoDB} para garantizar soporte ACID. Sobresale por su facilidad de despliegue, curva de aprendizaje inicial baja y alto rendimiento en operaciones de lectura simples.

\subsubsection*{MariaDB}
Bifurcación comunitaria de MySQL, compatible a nivel de binarios, pero con optimizadores de consultas más avanzados y un pool de hilos nativo para mejorar el manejo de la concurrencia.

\subsection*{Comparativa técnica}

En la Tabla \ref{tab:comparativa-bd} se resumen las principales características técnicas evaluadas.

\begin{table}[H]
    \hspace*{-0.25cm}
    \begin{tabular}{|p{0.22\textwidth}|p{0.22\textwidth}|p{0.22\textwidth}|p{0.22\textwidth}|}
        \hline
        \diagbox[width=\dimexpr 0.22\textwidth+2\tabcolsep\relax]{\textbf{Criterio}}{\textbf{Tecnología}} & \textbf{PostgreSQL} & \textbf{MySQL (InnoDB)} & \textbf{MariaDB} \\ \hline
        \textbf{Modelo} & Objeto-relacional & Relacional & Relacional \\ \hline
        \textbf{Cumplimiento SQL} & Estricto & Medio-alto & Medio-alto \\ \hline
        \textbf{Consultas complejas} & Excelente (CTE, sub consultas) & Adecuado & Bueno \\ \hline
        \textbf{Tipos avanzados} & Arrays, \texttt{JSONB} indexable, UUID & JSON básico & JSON básico \\ \hline
    \end{tabular}
    \caption{Comparativa de motores de bases de datos relacionales.}
    \label{tab:comparativa-bd}
\end{table}

\subsection*{Justificación de la elección}

El modelo de datos del dominio deportivo presenta una fuerte necesidad de integridad referencial, lo que descarta arquitecturas no relacionales. Se ha seleccionado \textbf{PostgreSQL} por los siguientes motivos:
\begin{itemize}
    \item \textbf{Integridad de datos:} Permite aplicar restricciones complejas directamente en la base de datos, asegurando la consistencia del dominio.
    \item \textbf{Flexibilidad estructural:} El uso de \texttt{JSONB} permite almacenar metadatos variables sin perder la rigidez relacional del resto del esquema.
    \item \textbf{Robustez a largo plazo:} Garantiza un rendimiento superior en consultas analíticas complejas (como la generación de clasificaciones o estadísticas).
\end{itemize}

Para gestionar la evolución del esquema y asegurar la reproducibilidad del entorno, se ha integrado \emph{Flyway}, herramienta que permite aplicar un control de versiones estricto sobre las migraciones \emph{.sql}.

Cabe comentar que para proyectos desarrollados con Spring Boot las alternativas principales son \emph{Flyway} y \emph{Liquibase}. Se ha optado por la primera y no por la segunda, ya que \emph{Liquibase} utiliza formatos como \texttt{XML}, \texttt{YAML} o \texttt{JSON} para definir los cambios, abstrayendo el motor de la base de datos y añadiendo una curva de aprendizaje. Esto sería positivo si se planease cambiar el motor de la base de datos en un futuro, pero para una implementación relativamente sencilla \emph{Flyway} es la alternativa óptima.

\section{Almacenamiento de objetos}

Para la gestión de archivos no estructurados generados por el sistema, como imágenes de perfil o documentos PDF (actas deportivas), se requiere un servicio de almacenamiento de objetos independiente de la base de datos para no degradar el rendimiento transaccional.

\subsection*{Evaluación de alternativas}

\subsubsection*{MinIO}
Servidor de almacenamiento de objetos de alto rendimiento, 100\% de código abierto, diseñado para infraestructuras locales o privadas. Es totalmente compatible con la API de Amazon S3, permitiendo una integración estandarizada.

\subsubsection*{Amazon S3}
El estándar de facto en la industria para el almacenamiento de objetos en la nube pública. Ofrece una disponibilidad hasta del 99.999999999\% y escalabilidad ilimitada, delegando la gestión de la infraestructura al proveedor.

\subsubsection*{Firebase Storage}
Servicio gestionado de Google orientado a integraciones rápidas en aplicaciones móviles y web, fuertemente acoplado al ecosistema de Firebase y a sus reglas de seguridad.

\subsection*{Comparativa técnica}

En la Tabla \ref{tab:comparativa-os} se analizan las soluciones basándonos en las necesidades del entorno de desarrollo y producción del proyecto.

\begin{table}[H]
    \hspace*{-0.25cm}
    \begin{tabular}{|p{0.22\textwidth}|p{0.22\textwidth}|p{0.22\textwidth}|p{0.22\textwidth}|}
        \hline
        \diagbox[width=\dimexpr 0.22\textwidth+2\tabcolsep\relax]{\textbf{Criterio}}{\textbf{Tecnología}} & \textbf{MinIO} & \textbf{Amazon S3} & \textbf{Firebase Storage} \\ \hline
        \textbf{Despliegue} & Local / Autoalojado (Docker) & Nube gestionada (AWS) & Nube gestionada (GCP) \\ \hline
        \textbf{Compatibilidad API} & Estándar S3 & Estándar S3 & Propietaria \\ \hline
        \textbf{Soberanía de datos} & Total (Bajo control propio) & Cedida al proveedor & Cedida al proveedor \\ \hline
        \textbf{Gestión de Infraestructura} & Responsabilidad del desarrollador & Automática & Automática \\ \hline
    \end{tabular}
    \caption{Comparativa de soluciones de almacenamiento de objetos.}
    \label{tab:comparativa-os}
\end{table}

\subsection*{Justificación de la elección}

Se ha seleccionado \textbf{MinIO} como motor de almacenamiento de objetos por las siguientes razones estratégicas:
\begin{itemize}
    \item \textbf{Soberanía y control:} Al desplegarse mediante contenedores de forma local, garantiza la privacidad total de los documentos y evita la dependencia (vendor lock-in) frente a los cambios de políticas de precios de terceros.
    \item \textbf{Entorno de desarrollo aislado:} Permite a cualquier desarrollador levantar el ecosistema completo sin necesidad de configurar credenciales en la nube ni incurrir en costes operativos durante la fase de creación.
    \item \textbf{Compatibilidad futura:} Al implementar exactamente la misma API que Amazon S3, el código desarrollado en el back-end es completamente agnóstico. Si en el futuro el proyecto requiere escalar a la nube pública, la migración a AWS S3 se podría realizar modificando únicamente variables de entorno, sin reescribir la lógica de la aplicación.
\end{itemize}

\section{Aplicación front-end}

Para la interfaz de usuario, se requiere una solución multiplataforma que evite la duplicidad de bases de código para web, iOS y Android, garantizando un rendimiento fluido y un acceso estandarizado a las funcionalidades nativas del dispositivo.

\subsection*{Evaluación de alternativas}

\subsubsection*{React Native con Expo}
Framework que renderiza componentes nativos reales comunicándose a través de un puente con el hilo de JavaScript/TypeScript. Expo abstrae la configuración de los entornos nativos, ofreciendo empaquetado rápido, recarga en caliente y gestión simplificada de despliegues.

\subsubsection*{Flutter}
Framework de Google basado en el lenguaje Dart que dibuja su propia interfaz mediante el motor gráfico Skia/Impeller. Ofrece un rendimiento visual sobresaliente en animaciones complejas, pero requiere adoptar un ecosistema completamente distinto.

\subsubsection*{Ionic / Capacitor}
Solución híbrida que ejecuta una aplicación web estándar (HTML/CSS/JS) dentro de un navegador interno (\emph{WebView}) embebido en la aplicación nativa. Destaca por su facilidad de desarrollo web, pero sufre limitaciones de rendimiento en interfaces pesadas.

\subsection*{Comparativa técnica}

En la Tabla \ref{tab:comparativa-frontend} se detalla la comparativa técnica de las tecnologías evaluadas.

\begin{table}[H]
    \hspace*{-0.25cm}
    \begin{tabular}{|p{0.22\textwidth}|p{0.22\textwidth}|p{0.22\textwidth}|p{0.22\textwidth}|}
        \hline
        \diagbox[width=\dimexpr 0.22\textwidth+2\tabcolsep\relax]{\textbf{Criterio}}{\textbf{Tecnología}} & \textbf{React Native + Expo} & \textbf{Flutter} & \textbf{Ionic} \\ \hline
        \textbf{Renderizado UI} & Componentes Nativos & Motor gráfico propio & WebView \\ \hline
        \textbf{Lenguaje} & TypeScript / JavaScript & Dart & TypeScript / Web \\ \hline
        \textbf{Acceso a Hardware} & Mediante puente (JSI) & Integración compilada & Plug-ins (Capacitor) \\ \hline
        \textbf{Desarrollo ágil} & Muy alto (Expo) & Alto & Muy alto \\ \hline
    \end{tabular}
    \caption{Comparativa de frameworks para el desarrollo cliente.}
    \label{tab:comparativa-frontend}
\end{table}

\subsection*{Justificación de la elección}

Se ha seleccionado \textbf{React Native junto con Expo} frente a sus alternativas por las siguientes razones:
\begin{itemize}
    \item \textbf{Experiencia Nativa:} Se descarta Ionic para evitar las limitaciones de rendimiento táctil y visual que imponen los \emph{WebViews}, buscando ofrecer una experiencia de usuario genuinamente nativa.
    \item \textbf{Reutilización de paradigmas:} Al poseer experiencia previa en React, se elimina la curva de aprendizaje de un nuevo lenguaje como Dart. Esto permite invertir el tiempo del proyecto en resolver la lógica de negocio y optimizar las peticiones a la API.
    \item \textbf{Aceleración del ciclo de vida:} Las herramientas de empaquetado de Expo facilitan las pruebas continuas en dispositivos físicos mediante códigos QR, sin necesidad de lidiar con las configuraciones locales de Xcode o Android Studio durante el desarrollo.
\end{itemize}
